\subsubsection{Kernel Logistic Regression}

\paragraph{Ý tưởng}
Kernel Logistic Regression là mở rộng của Logistic Regression nhằm
xử lý các bài toán phân loại phi tuyến thông qua \textit{kernel trick}.
Thay vì học một vector trọng số trong không gian đặc trưng ban đầu,
mô hình được biểu diễn trong dạng đối ngẫu (dual form), trong đó hàm
quyết định phụ thuộc vào các điểm huấn luyện thông qua một hàm kernel.

Hàm quyết định của mô hình được viết dưới dạng:
\begin{equation}
  f(x) = \sum_{i=1}^{n} \alpha_i t_i K(\boldsymbol{x}_i, \boldsymbol{x}),
\end{equation}
trong đó $\alpha_i$ là các tham số cần học, $t_i \in \{-1, 1\}$ là nhãn
huấn luyện, và $K(\boldsymbol{x}_i, x)$ là hàm kernel.

Trong trường hợp sử dụng kernel Gaussian RBF, ta có:
\begin{equation}
  K(\boldsymbol{x}, \boldsymbol{x'}) = \exp\left(-\gamma \|\boldsymbol{x} -
  \boldsymbol{x}'\|^2\right).
\end{equation}

Xác suất dự đoán được mô hình hóa thông qua hàm sigmoid:
\begin{equation}
  P(t=1 \mid x) = \sigma(f(\boldsymbol{x}) = \frac{1}{1 +
  \exp(-f(\boldsymbol{x}))}.
\end{equation}

Hàm mất mát của Kernel Logistic Regression là log-loss trong không
gian kernel, kết hợp với regularization:
\begin{equation}
  \mathcal{L}(\boldsymbol{\alpha})
  = \sum_{i=1}^{n}
  \log\left(1 + \exp\left(-t_i f(\boldsymbol{x}_i)\right)\right)
  + \lambda \, \boldsymbol{\alpha}^\intercal \boldsymbol{K}
  \boldsymbol{\alpha},
\end{equation}
trong đó $\boldsymbol{K}$ là ma trận kernel với
$K_{ij} = K(\boldsymbol{x}_i, \boldsymbol{x}_j)$.

Việc tối ưu mô hình được thực hiện trong không gian đối ngẫu, thay vì
tối ưu trực tiếp trên vector trọng số như Logistic Regression tuyến
tính. Điều này cho phép mô hình học được các ranh giới quyết định phi
tuyến thông qua ánh xạ ngầm sang không gian đặc trưng có chiều cao
(thậm chí vô hạn chiều).

\paragraph{So sánh Logistic Regression và Kernel Logistic Regression}
Logistic Regression là mô hình tuyến tính, trong đó hàm quyết định
được biểu diễn dưới dạng:
\begin{equation}
  f(\boldsymbol{x}) = \boldsymbol{w}^\intercal \boldsymbol{x} + b.
\end{equation}

Do đó, ranh giới quyết định của mô hình luôn là một siêu phẳng trong
không gian đầu vào. Điều này khiến Logistic Regression phù hợp với
các bài toán có cấu trúc tuyến tính hoặc gần tuyến tính, nhưng gặp
hạn chế lớn khi dữ liệu có phân bố phi tuyến.

Ngược lại, Kernel Logistic Regression thay thế tích vô hướng trực tiếp
bằng hàm kernel:
\begin{equation}
  f(\boldsymbol{x}) = \sum_{i=1}^{n} \alpha_i t_i K(\boldsymbol{x}_i,
  \boldsymbol{x}).
\end{equation}

Nhờ kernel trick, mô hình có thể ánh xạ ngầm dữ liệu sang không gian
đặc trưng có chiều cao hơn, giúp biểu diễn được các quan hệ phi tuyến
phức tạp mà không cần tính toán tường minh ánh xạ.

Về mặt hình học, Logistic Regression tạo ra ranh giới quyết định
tuyến tính, trong khi Kernel Logistic Regression có thể tạo ra ranh
giới cong linh hoạt tùy thuộc vào lựa chọn kernel (ví dụ RBF kernel).

Về độ phức tạp, Logistic Regression có chi phí tính toán thấp hơn và
mở rộng tốt với dữ liệu lớn. Ngược lại, Kernel Logistic Regression
phụ thuộc vào ma trận kernel kích thước $n \times n$, dẫn đến chi phí
tính toán và bộ nhớ cao hơn.

Tóm lại, Logistic Regression phù hợp với dữ liệu tuyến tính và bài
toán lớn, trong khi Kernel Logistic Regression phù hợp với dữ liệu
phi tuyến nhưng đánh đổi về chi phí tính toán.

\paragraph{So sánh trên dữ liệu không phân chia tuyến tính}
\figurename~\ref{fig:logistic-regression-xor} và
\figurename~\ref{fig:kernel-logistic-regression-xor} thể hiện đường biên quyết
định của Logistic Regression và Kernel Logistic Regression trên dữ liệu không
phân chia tuyến tính. Dễ thấy với Logistic Regression thì đường biên này là
một đường thẳng do đó vẫn còn nhiều điểm không thể phân loại đúng. Trong khi
đó Kernel Logistic Regression có đường biên quyết định là đường cong phi tuyến
do đó có thể phân loại được những điểm này và gần như phân loại đúng tất cả
trừ vài điểm dữ liệu nhiễu.

\begin{figure}[htbp]
  \centering
  \includegraphics[width=0.9\textwidth]{logistic-regression-xor}
  \caption{Đường biên quyết định của Logistic Regression cho dữ liệu không
  phân chia tuyến tính}
  \label{fig:logistic-regression-xor}
\end{figure}

\begin{figure}[htbp]
  \centering
  \includegraphics[width=0.9\textwidth]{kernel-logistic-regression-xor}
  \caption{Đường biên quyết định của Kernel Logistic Regression cho dữ liệu
  không phân chia tuyến tính}
  \label{fig:kernel-logistic-regression-xor}
\end{figure}

\paragraph{Kết luận}
Từ kết quả thực nghiệm, có thể thấy rằng Logistic Regression chỉ
phù hợp với các bài toán có ranh giới quyết định tuyến tính, do
giới hạn trong việc học một siêu phẳng trong không gian đầu vào.

Ngược lại, Kernel Logistic Regression có khả năng mô hình hóa các
quan hệ phi tuyến nhờ kernel trick, cho phép xây dựng ranh giới
quyết định linh hoạt hơn và phù hợp với các cấu trúc dữ liệu phức
tạp như XOR.

Tuy nhiên, khả năng biểu diễn mạnh hơn cũng đi kèm với chi phí tính
toán cao hơn và độ nhạy với tham số kernel, đặc biệt là trong các
trường hợp dữ liệu có nhiễu hoặc số lượng mẫu lớn.

Do đó, việc lựa chọn giữa hai mô hình cần cân bằng giữa độ phức tạp
của dữ liệu và yêu cầu về hiệu năng tính toán.
